% !TeX root = part5.tex
\documentclass[../final.tex]{subfiles}

\begin{document}

\practice{5}{}{Рентгеновские спектры. Гелиеподобные атомы.
Теория возмущений и вариационный метод}
\captionsetup{font=footnotesize,labelfont=bf,skip=5pt}

\begingroup
\renewcommand{\thesection}{2.\arabic{section}}

\task[6]{Поправка экранирования для линии \texorpdfstring{$K_\alpha$}{K-alpha}}

\condition
Определить поправку экранирования $\sigma_K$ в законе Мозли
для $K_\alpha$-линий Fe, Co, Sn и Cs. Длины волн соответственно равны
\[
    1.935\,\text{\AA},\quad 1.787\,\text{\AA},\quad
    0.492\,\text{\AA},\quad 0.402\,\text{\AA}.
\]

\solution
В \rassumption{нерелятивистской модели с эффективным зарядом}
$Z_{\mathrm{eff}}=Z-\sigma_K$ переход $K_\alpha$ соответствует
переходу из оболочки $L$ в оболочку $K$, то есть $n=2\to n=1$:
\[
    E_\gamma=R_E(Z-\sigma_K)^2
        \left(\frac1{1^2}-\frac1{2^2}\right)
        =\frac34R_E(Z-\sigma_K)^2.
\]
Здесь $R_E=13.606\,\text{эВ}$ --- энергетическая постоянная Ридберга.
Следовательно,
\[
    \sigma_K=Z-\sqrt{\frac{4E_\gamma}{3R_E}}
        =Z-\sqrt{\frac{4hc}{3R_E\lambda}},
    \qquad hc=12\,398.42\,\text{эВ}\cdot\text{\AA}.
\]
Например, для железа
\[
    E_\gamma=\frac{12\,398.42}{1.935}\,\text{эВ}
       \approx6.407\,\text{кэВ},\qquad
    \sigma_K=26-\sqrt{\frac{4\cdot6407.45}{3\cdot13.606}}
       \approx0.942.
\]

\answer
\begin{center}
\begin{tabular}{c|r|c|c|c}
    \rconcept{Атом} & $Z$ & $\lambda$, \text{\AA}
        & $E_\gamma$, кэВ & $\sigma_K$\\ \hline
    Fe & 26 & 1.935 & 6.407 & 0.942\\
    Co & 27 & 1.787 & 6.938 & 0.925\\
    Sn & 50 & 0.492 & 25.200 & 0.305\\
    Cs & 55 & 0.402 & 30.842 & 0.023
\end{tabular}
\end{center}
Эта эффективная поправка определяется по выбранной модели спектра;
она не обязана быть одинаковой для разных элементов.

\task[7]{Фотоэлектрон из \texorpdfstring{$K$}{K}-оболочки молибдена}

\condition
Найти кинетическую энергию электрона, вырываемого из $K$-оболочки
молибдена $K_\alpha$-излучением серебра. Используем
\[
    \lambda_{K_\alpha}(\mathrm{Ag})=0.56\,\text{\AA},\qquad
    E_{\text{св}}(K,\mathrm{Mo})\approx20\,\text{кэВ}.
\]

\solution
Пренебрегая отдачей атома, запишем баланс энергии при фотоэффекте:
\[
    h\nu=E_{\text{св}}+T.
\]
Энергия фотона серебра равна
\[
    h\nu=\frac{hc}{\lambda_{K_\alpha}}
       =\frac{12\,398.42}{0.56}\,\text{эВ}
       \approx22.14\,\text{кэВ}.
\]
Энергия связи $K$-электрона характеризует весь атом молибдена,
а не изолированный водородоподобный ион. В частности, её нельзя
без дополнительной модели заменить энергией $1s$-уровня при $Z=42$.

\answer
\[
    T\approx22.14-20=2.14\,\text{кэВ}\approx2.1\,\text{кэВ}.
\]

\endgroup

\rtheory{Два электрона: спин и принцип Паули}

\subsection*{Гамильтониан и единицы}

Рассматриваем два электрона в поле неподвижного ядра с зарядом $Ze$.
В \rassumption{СГС}, без релятивистских и спин-зависимых поправок,
\begin{align*}
    \hat H&=\hat H_0+\hat V,\qquad
    \hat H_0=\hat h_1+\hat h_2,\\
    \hat h_i&=-\frac{\hbar^2}{2m_e}\nabla_i^2-\frac{Ze^2}{r_i},
    \qquad \hat V=\frac{e^2}{r_{12}},
    \qquad r_{12}=|\mathbf r_1-\mathbf r_2|.
\end{align*}
В СИ вместо $e^2$ следует писать $e^2/(4\pi\varepsilon_0)$.
Для расчётов введём боровский радиус и атомную единицу энергии:
\[
    a_0=\frac{\hbar^2}{m_e e^2},\qquad
    E_{\mathrm h}=\frac{e^2}{a_0}=2R_E\approx27.211\,\text{эВ}.
\]
Энергии отсчитываем от состояния свободного ядра и двух свободных
электронов: энергия связанного атома отрицательна.

\subsection*{Спиновые функции}

Для одного электрона
\[
    \alpha=\begin{pmatrix}1\\0\end{pmatrix},\quad m_s=\frac12,
    \qquad
    \beta=\begin{pmatrix}0\\1\end{pmatrix},\quad m_s=-\frac12.
\]
\rconcept{Парагелий} имеет полный спин $S=0$.
Его спиновая функция антисимметрична:
\[
    \chi_{00}=\frac{\alpha(1)\beta(2)-\beta(1)\alpha(2)}{\sqrt2}.
\]
\rconcept{Ортогелий} имеет $S=1$. Три спиновые функции симметричны:
\begin{align*}
    \chi_{11}&=\alpha(1)\alpha(2),\\
    \chi_{10}&=\frac{\alpha(1)\beta(2)+\beta(1)\alpha(2)}{\sqrt2},\\
    \chi_{1,-1}&=\beta(1)\beta(2).
\end{align*}

\subsection*{Пространственная симметрия}

Электроны --- тождественные фермионы. По \rconcept{принципу Паули}
полная функция при перестановке электронов меняет знак:
\[
    \hat P_{12}\Psi=-\Psi.
\]
Для двух различных ортонормированных орбиталей $a$ и $b$ обозначим
\[
    \Phi_1=\varphi_a(\mathbf r_1)\varphi_b(\mathbf r_2),\qquad
    \Phi_2=\varphi_b(\mathbf r_1)\varphi_a(\mathbf r_2),
    \qquad \Phi_\pm=\frac{\Phi_1\pm\Phi_2}{\sqrt2}.
\]
Тогда допустимые полные функции имеют вид
\[
    \Psi_{\text{пара}}=\Phi_+\chi_{00},\qquad
    \Psi_{\text{орто},M_S}=\Phi_-\chi_{1M_S},
    \qquad M_S=1,0,-1.
\]
Синглету соответствует симметричная пространственная функция,
триплету --- антисимметричная.
Если оба электрона занимают одну орбиталь, то антисимметричная
пространственная комбинация равна нулю. Поэтому конфигурация $1s^2$
допускает только терм ${}^1S_0$, но не ${}^3S_1$.

\begingroup
\renewcommand{\thesection}{3.\arabic{section}}

\task[1]{Основное состояние: первый порядок теории возмущений}
\label{task:fiham5_ground_perturbation}

\condition
Найти поправку к энергии основного состояния гелиеподобного иона
в первом порядке теории возмущений по межэлектронному отталкиванию.

\solution
В нулевом приближении оба электрона находятся на водородоподобной
$1s$-орбитали:
\[
    \varphi_{1s}(\mathbf r)=\frac1{\sqrt\pi}
        \left(\frac Z{a_0}\right)^{3/2}e^{-Zr/a_0},\qquad
    \Phi_0=\varphi_{1s}(\mathbf r_1)\varphi_{1s}(\mathbf r_2).
\]
Полная функция равна $\Psi_0=\Phi_0\chi_{00}$, а энергия без
взаимодействия электронов составляет
\[
    E_0^{(0)}=-2Z^2R_E.
\]
Поскольку $\hat V$ не действует на спин, поправка вычисляется
по пространственной функции:
\[
    E_0^{(1)}=\langle\Phi_0|\hat V|\Phi_0\rangle
        =e^2\iint\frac{|\Phi_0|^2}{r_{12}}\,d^3r_1\,d^3r_2.
\]

\subsection*{Вычисление кулоновского интеграла}

Используем безразмерный радиус $\rho=r/a_0$.
Радиальная плотность вероятности для $1s$-электрона равна
\[
    p_Z(\rho)=4Z^3\rho^2e^{-2Z\rho},\qquad
    \int_0^\infty p_Z(\rho)\,d\rho=1.
\]
Угловое усреднение $1/|\boldsymbol\rho_1-\boldsymbol\rho_2|$
для сферически симметричных плотностей даёт $1/\max(\rho_1,\rho_2)$.
Разбивая область интегрирования на две симметричные части, получаем
\[
    \frac{E_0^{(1)}}{E_{\mathrm h}}
      =2\int_0^\infty\frac{p_Z(\rho)}\rho
        \left[\int_0^\rho p_Z(s)\,ds\right]d\rho.
\]
Внутренний интеграл равен
\[
    \int_0^\rho p_Z(s)\,ds
      =1-e^{-2Z\rho}\bigl(1+2Z\rho+2Z^2\rho^2\bigr).
\]
После подстановки и интегрирования
\begin{align*}
    \frac{E_0^{(1)}}{E_{\mathrm h}}
      &=8Z^3\int_0^\infty\rho e^{-2Z\rho}
        \left[1-e^{-2Z\rho}
        \bigl(1+2Z\rho+2Z^2\rho^2\bigr)\right]d\rho\\
      &=\frac58Z.
\end{align*}

\answer
\[
    E_0^{(1)}=\frac58ZE_{\mathrm h}=\frac54ZR_E,\qquad
    E_0\approx\left(-2Z^2+\frac54Z\right)R_E.
\]
Для гелия, $Z=2$,
\[
    E_0^{(0)}\approx-108.85\,\text{эВ},\qquad
    E_0^{(1)}\approx34.01\,\text{эВ},\qquad
    E_0\approx-74.83\,\text{эВ}.
\]
Отталкивание электронов \rresult{повышает полную энергию}.

\task[2]{Вариационный расчёт энергий ионизации He и Li\texorpdfstring{$^+$}{+}}
\label{task:fiham5_variational}

\condition
Вариационным методом определить энергии ионизации атома гелия
и иона лития $\mathrm{Li}^+$. В качестве пробных функций использовать
водородоподобные $1s$-орбитали с эффективным зарядом.

\solution
Выберем нормированную пробную пространственную функцию
\[
    \Phi_\zeta(\mathbf r_1,\mathbf r_2)
       =\varphi_\zeta(\mathbf r_1)\varphi_\zeta(\mathbf r_2),\qquad
    \varphi_\zeta(\mathbf r)=\frac1{\sqrt\pi}
        \left(\frac\zeta{a_0}\right)^{3/2}e^{-\zeta r/a_0},
    \quad \zeta=Z_{\mathrm{eff}}>0.
\]
Спиновая часть --- $\chi_{00}$. В гамильтониане остаётся настоящий
заряд ядра $Z$, тогда как $\zeta$ варьируется только в пробной функции.

\subsection*{Средние кинетической и потенциальной энергий}

Положим $\alpha=\zeta/a_0$ и
$\varphi_\zeta=\mathcal N e^{-\alpha r}$.
Для сферически симметричной функции
\[
    \nabla^2\varphi_\zeta
       =\frac1{r^2}\frac{d}{dr}
          \left(r^2\frac{d\varphi_\zeta}{dr}\right)
       =\left(\alpha^2-\frac{2\alpha}r\right)\varphi_\zeta,
    \qquad \left\langle\frac1r\right\rangle=\alpha.
\]
Следовательно, для каждого электрона
\begin{align*}
    \langle\hat T_i\rangle
       &=-\frac{\hbar^2}{2m_e}
           \left(\alpha^2-2\alpha\left\langle\frac1r\right\rangle\right)
         =\zeta^2R_E,\\
    \left\langle-\frac{Ze^2}{r_i}\right\rangle
       &=-\frac{Z\zeta e^2}{a_0}=-2Z\zeta R_E.
\end{align*}
Из кулоновского интеграла предыдущей задачи с заменой $Z\to\zeta$
\[
    \langle\hat V\rangle=\frac54\zeta R_E.
\]
Таким образом,
\[
    E(\zeta)=\left(2\zeta^2-4Z\zeta+\frac54\zeta\right)R_E.
\]

\subsection*{Минимизация и экранирование}

\begin{align*}
    \frac{dE}{d\zeta}
       &=\left(4\zeta-4Z+\frac54\right)R_E=0,\\
    \zeta_*&=Z-\frac5{16},\qquad
    \frac{d^2E}{d\zeta^2}=4R_E>0,\\
    E_{\mathrm{var}}&=-2\left(Z-\frac5{16}\right)^2R_E.
\end{align*}
В этой модели второй электрон уменьшает эффективный заряд
на $5/16$. Согласно вариационному принципу полученная энергия
является верхней оценкой точной энергии основного состояния.

\subsection*{Энергия ионизации}

После удаления одного электрона остаётся водородоподобный ион
с энергией $E_{\text{ион}}=-Z^2R_E$.
Минимальная энергия удаления электрона равна
\[
    I=E_{\text{ион}}-E_{\mathrm{var}}
       =\left[2\left(Z-\frac5{16}\right)^2-Z^2\right]R_E.
\]

\answer
\begin{center}
\begin{tabular}{c|c|c|c|c}
    \rconcept{Система} & $Z$ & $Z_{\mathrm{eff}}$
        & $E_{\mathrm{var}}$, эВ & $I$, эВ\\ \hline
    He & 2 & $27/16$ & $-77.49$ & $23.07$\\
    $\mathrm{Li}^+$ & 3 & $43/16$ & $-196.54$ & $74.09$
\end{tabular}
\end{center}
Значения относятся к \rassumption{однопараметрической пробной функции}.

\task[3]{Минимальная энергия ортогелия и обменное расщепление}
\label{task:fiham5_ortho}

\condition
Теорией возмущений определить минимальную полную электронную энергию
ортогелия, ограничившись водородоподобными орбиталями $1s$ и $2s$.
Найти в этом же приближении энергию возбуждения из основного состояния
парагелия в ортогелий.

\solution

\subsection*{Конфигурация и невозмущённая энергия}

Для триплета пространственная функция должна быть антисимметрична.
Конфигурация $1s^2$ невозможна, поскольку $\Phi_-=0$.
В выбранном наборе орбиталей минимальная конфигурация ортогелия ---
$1s2s$. Обе орбитали имеют $l=0$, поэтому $L=0$, $S=1$, $J=1$:
\[
    1s2s\;{}^3S_1,
    \qquad \text{обозначение терма: }{}^{2S+1}L_J.
\]
Невозмущённые одноэлектронные уровни равны
$E_n^{(0)}=-Z^2R_E/n^2$, откуда
\[
    E_{1s2s}^{(0)}=-Z^2R_E-\frac14Z^2R_E
       =-\frac54Z^2R_E.
\]
При $Z=2$ получаем $E_{1s2s}^{(0)}=-5R_E\approx-68.03\,\text{эВ}$.

\subsection*{Матрица возмущения}

В базисе $\Phi_1=\varphi_{1s}(1)\varphi_{2s}(2)$,
$\Phi_2=\varphi_{2s}(1)\varphi_{1s}(2)$ матрица взаимодействия имеет вид
\[
    V=\begin{pmatrix}C&A\\A&C\end{pmatrix}.
\]
Здесь \rconcept{кулоновский интеграл}
\[
    C=\iint |\varphi_{1s}(\mathbf r_1)|^2
       \frac{e^2}{r_{12}}|\varphi_{2s}(\mathbf r_2)|^2
       \,d^3r_1\,d^3r_2,
\]
а \rconcept{обменный интеграл}
\[
    A=\iint\varphi_{1s}^*(\mathbf r_1)\varphi_{2s}^*(\mathbf r_2)
       \frac{e^2}{r_{12}}
       \varphi_{2s}(\mathbf r_1)\varphi_{1s}(\mathbf r_2)
       \,d^3r_1\,d^3r_2.
\]
Собственные пространственные функции --- $\Phi_+$ и $\Phi_-$:
\begin{align*}
    E_{\text{пара}}&=E_{1s2s}^{(0)}+C+A,
        &&S=0,\quad 1s2s\;{}^1S_0,\\
    E_{\text{орто}}&=E_{1s2s}^{(0)}+C-A,
        &&S=1,\quad 1s2s\;{}^3S_1.
\end{align*}
Таким образом, $E_{\text{пара}}-E_{\text{орто}}=2A$.
При $A>0$ триплет лежит ниже синглета той же конфигурации.

\subsection*{Одноэлектронные функции}

В центральном поле
\[
    \varphi_{nlm}(\mathbf r)=R_{nl}(r)Y_{lm}(\Omega),\qquad
    Y_{00}=\frac1{\sqrt{4\pi}}.
\]
Используем безразмерные координаты $\boldsymbol\rho=\mathbf r/a_0$
и нормированные безразмерные функции
$\widetilde\varphi=a_0^{3/2}\varphi$:
\begin{align*}
    \widetilde\varphi_{1s}(\rho)
       &=\frac{Z^{3/2}}{\sqrt\pi}e^{-Z\rho},\\
    \widetilde\varphi_{2s}(\rho)
       &=\frac{Z^{3/2}}{4\sqrt{2\pi}}
           (2-Z\rho)e^{-Z\rho/2}.
\end{align*}

\subsection*{Угловое интегрирование}

Пусть $\rho_< =\min(\rho_1,\rho_2)$,
$\rho_> =\max(\rho_1,\rho_2)$, а $\theta$ --- угол между радиус-векторами.
Из производящей функции полиномов Лежандра
\[
    (1-2tx+t^2)^{-1/2}=\sum_{\ell=0}^\infty t^\ell P_\ell(x),
    \qquad |t|<1,
\]
получаем разложение
\[
    \frac1{\rho_{12}}
       =\sum_{\ell=0}^\infty
          \frac{\rho_<^{\ell}}{\rho_>^{\ell+1}}P_\ell(\cos\theta).
\]
По теореме сложения сферических гармоник
\[
    P_\ell(\cos\theta)=\frac{4\pi}{2\ell+1}
       \sum_{m=-\ell}^{\ell}Y_{\ell m}^*(\Omega_1)Y_{\ell m}(\Omega_2).
\]
Обе орбитали сферически симметричны, поэтому после углового
интегрирования остаётся только член $\ell=0$:
\[
    \int Y_{00}\,d\Omega=\sqrt{4\pi},\qquad
    \iint\frac{d\Omega_1\,d\Omega_2}{\rho_{12}}
       =\frac{16\pi^2}{\rho_>}.
\]
Радиальное ядро имеет вид
\[
    \frac1{\rho_>}=
       \begin{cases}
           1/\rho_1,&\rho_2<\rho_1,\\
           1/\rho_2,&\rho_2>\rho_1.
       \end{cases}
\]

\subsection*{Кулоновский интеграл}

Учитывая $d^3\rho=\rho^2d\rho\,d\Omega$, получаем
\begin{align*}
    \frac C{E_{\mathrm h}}
      =\frac{Z^6}{2}\int_0^\infty d\rho\,\rho^2e^{-2Z\rho}
       \Biggl[&\frac1\rho\int_0^\rho
          ds\,s^2(2-Zs)^2e^{-Zs}\\
          &+\int_\rho^\infty ds\,s(2-Zs)^2e^{-Zs}\Biggr].
\end{align*}
Для степеней радиуса используем
\[
    \int_0^\infty x^n e^{-ax}\,dx=\frac{n!}{a^{n+1}},\qquad a>0.
\]
После вычисления радиальных интегралов
\[
    C=\frac{17}{81}ZE_{\mathrm h}=\frac{34}{81}ZR_E.
\]

\subsection*{Обменный интеграл}

Произведение функций в одной точке равно
\[
    \widetilde\varphi_{1s}(\rho)\widetilde\varphi_{2s}(\rho)
       =\frac{Z^3}{4\sqrt2\,\pi}(2-Z\rho)e^{-3Z\rho/2}.
\]
Радиальная подынтегральная функция обменного интеграла симметрична
по $\rho_1$ и $\rho_2$. Поэтому достаточно удвоить интеграл по области
$\rho_2<\rho_1$:
\[
    \frac A{E_{\mathrm h}}
      =Z^6\int_0^\infty d\rho\,\rho(2-Z\rho)e^{-3Z\rho/2}
         \int_0^\rho ds\,s^2(2-Zs)e^{-3Zs/2}.
\]
Тем же способом получаем
\[
    A=\frac{16}{729}ZE_{\mathrm h}=\frac{32}{729}ZR_E.
\]
При $Z=2$
\[
    C\approx11.42\,\text{эВ},\qquad A\approx1.19\,\text{эВ}.
\]

\subsection*{Энергии термов и возбуждения}

Полная энергия низшего триплета в выбранном приближении равна
\begin{align*}
    E_{\text{орто}}
       &=\left[-\frac54Z^2+
             \left(\frac{34}{81}-\frac{32}{729}\right)Z\right]R_E\\
       &=\left(-\frac54Z^2+\frac{274}{729}Z\right)R_E.
\end{align*}
Для гелия
\[
    E_{\text{орто}}\approx-57.80\,\text{эВ},\qquad
    E_{\text{пара}}(1s2s)\approx-55.41\,\text{эВ},\qquad
    2A\approx2.39\,\text{эВ}.
\]
Энергию основного состояния $1s^2\;{}^1S_0$ берём
из задачи~\ref{task:fiham5_ground_perturbation}, то есть также
в первом порядке по $e^2/r_{12}$:
\[
    E_{\text{осн}}=\left(-2Z^2+\frac54Z\right)R_E,
    \qquad E_{\text{осн}}(Z=2)\approx-74.83\,\text{эВ}.
\]
Отсюда энергия возбуждения в ортогелий
\begin{align*}
    \Delta E&=E_{\text{орто}}-E_{\text{осн}}\\
       &=\left[\frac34Z^2+
          \left(\frac{274}{729}-\frac54\right)Z\right]R_E,
    \qquad \Delta E(Z=2)\approx17.03\,\text{эВ}.
\end{align*}

\begin{rbox}[breakable=false,left=16pt,right=16pt,top=14pt,bottom=10pt]{[]}
    \centering
    \begin{tikzpicture}[x=1cm,y=0.7cm,
        every node/.style={text=rtext,font=\small}]
        \draw[rtext,-{Stealth[length=2.2mm]}] (0,-0.3)--(0,5.5) node[above] {$E$};
        \draw[rtext,line width=1pt] (0.6,0)--(8.7,0);
        \node[above] at (6.0,0) {$1s^2\;{}^1S_0$};
        \draw[p3,line width=1.2pt] (3.3,3.5)--(8.7,3.5);
        \node[p3,above] at (6.0,3.5) {$1s2s\;{}^3S_1$};
        \draw[red3,line width=1.2pt] (3.3,4.9)--(8.7,4.9);
        \node[red3,above] at (6.0,4.9) {$1s2s\;{}^1S_0$};
        \draw[p3,{Stealth[length=2mm]}-{Stealth[length=2mm]}]
            (2.0,0.1)--(2.0,3.4) node[midway,left] {$\Delta E$};
        \draw[rtext,{Stealth[length=2mm]}-{Stealth[length=2mm]}]
            (9.4,3.6)--(9.4,4.8) node[midway,right] {$2A$};
    \end{tikzpicture}
    \captionof{figure}{Основной синглет и термы конфигурации $1s2s$.
    Триплет ниже возбуждённого синглета на $2A$.
    Вертикальные интервалы условные; стрелки обозначают разности энергий.}
    \label{fig:fiham5_exchange}
\end{rbox}

\answer
В первом порядке теории возмущений с неэкранированными
водородоподобными функциями
\[
    E_{\text{орто}}^{\min}\approx-57.80\,\text{эВ},\qquad
    \Delta E\approx17.03\,\text{эВ}.
\]
Это \rassumption{оценки в заданном приближении}, а не точные
спектроскопические энергии.

\endgroup
\end{document}
